\section{Resultados III --- Comparação pré/pós-ChatGPT}\label{cap:comparacao}

Este é o capítulo central da revisão. Ele compara, dimensão a dimensão, os 240
estudos publicados até 2022 com os 576 publicados de 2023 em diante, conforme a
Tabela~\ref{tab:comparacao}. O período é definido de forma determinística pelo
ano de publicação (Seção~\ref{sec:periodo}); como o ChatGPT foi lançado em 30 de
novembro de 2022, o ano de 2023 é o primeiro ano-calendário posterior ao pivô. A
hipótese de trabalho sustenta que a literatura pré-ChatGPT convergia para a
polarização com risco em tarefas rotineiras de baixa qualificação, ao passo que
a literatura pós-ChatGPT deslocaria o risco para tarefas cognitivas de alta
qualificação. Antes de discutir os eixos, é preciso fixar a natureza dos testes:
o corpus é o \emph{censo} dos estudos incluídos, não uma amostra aleatória de
uma população; os valores-$p$ de associação ($\chi^2$ e Fisher exato) são, por
isso, lidos como medidas exploratórias de associação, e não como inferência
populacional. Todas as proporções excluem do denominador os estudos que não
classificaram a dimensão.

\begin{table}[htbp]
\centering
\caption{Comparação pareada de dimensões pré e pós-ChatGPT.}
\label{tab:comparacao}
\input{tables/comparacao_pre_pos.tex}
\end{table}

\subsection{Eixo 1: Quem está em risco}

O eixo que testa diretamente a hipótese é a polarização
(Tabela~\ref{tab:polarizacao_prepos}). A participação de estudos que situam o
risco nas ocupações de \emph{alta} qualificação multiplica-se por cinco entre os
períodos: de 2,7\% (\gls{ic} Wilson 95\%: [1,1; 6,8]) no pré para 13,6\%
([10,5; 17,4]) no pós, com razão de chances de 5,61 e $p < 0{,}001$ no teste
exato de Fisher. Antes de 2021 não há, no corpus, um único estudo que situe o
risco primordialmente na alta qualificação; em 2026 eles são quase um em cada
cinco. A direção é a prevista pela hipótese, e a associação aparece também na
distribuição \emph{completa} da polarização ($\chi^2 = 13{,}58$,
$p = 0{,}0035$). Ilustram esse enquadramento estudos pós-ChatGPT que situam o
risco em ocupações cognitivas qualificadas, como a análise da escolha
ocupacional diante da IA generativa de \textcite{gollerGschwendtWolter2025} ---
cujo próprio título, ``\textit{This time it's different}'', anuncia a tese da
ruptura --- e o exame das ocupações no centro da adoção de IA de
\textcite{fontanelli2025}.

Dois qualificadores, porém, delimitam o alcance do achado. Primeiro, a categoria
modal continua sendo o risco para a \emph{baixa} qualificação nos dois períodos
(63,3\% no pré, 53,7\% no pós) --- padrão que ecoa a literatura da automação
sobre trabalhadores rotineiros e desigualdade, a exemplo de
\textcite{mollRachelRestrepo2022} ---, e a categoria ``ambos'' mantém peso
elevado e estável (32,7\% e 31,4\%). O crescimento do risco atribuído à alta
qualificação se faz, portanto, à custa da hegemonia da baixa qualificação, mas
sem revogá-la. Segundo, quando se restringe a análise aos estudos de maior
qualidade metodológica (\textit{score} $\geq 4$,
Tabela~\ref{tab:robustez_qualidade}), a elevação persiste em direção (5,9\%
para 12,5\%), mas não se distingue do acaso ($p = 0{,}5128$), com apenas 34
estudos classificados no pré. A Tabela~\ref{tab:h1_sensibilidade}, discutida no
Capítulo~\ref{cap:discussao}, mostra que o resultado resiste a oito outras
especificações e que a perda de significância nesse subconjunto decorre mais do
tamanho das células do que de uma reversão do efeito. A leitura é de
\emph{ruptura parcial}: o enquadramento do risco para tarefas cognitivas de alta
qualificação cresce de forma nítida e robusta após o ChatGPT, mas permanece
minoritário e não desloca o diagnóstico de polarização herdado.

\begin{table}[htbp]
\centering
\caption{Foco da hipótese: risco para alta qualificação, pré vs.\ pós-ChatGPT.}
\label{tab:polarizacao_prepos}
\input{tables/polarizacao_pre_pos.tex}
\end{table}

\begin{table}[htbp]
\centering
\caption{Robustez: polarização restrita a estudos de alta qualidade (score $\geq$ 4).}
\label{tab:robustez_qualidade}
\input{tables/robustez_qualidade.tex}
\end{table}

\subsection{Eixo 2: Qual o mecanismo}

Três dos quatro mecanismos do arcabouço de Acemoglu e Restrepo mudam de peso
entre os períodos. A \emph{complementaridade} sobe de 57,8\% para 73,0\% dos
estudos ($p < 0{,}001$). O \emph{deslocamento} permanece o mecanismo mais
invocado em ambos os períodos, mas deixa de ser quase universal: recua de 95,4\%
para 85,1\% ($p < 0{,}001$). A \emph{demanda agregada} cai praticamente pela
metade, de 23,5\% para 12,4\% ($p < 0{,}001$). Apenas a \emph{reinstalação}
permanece estável (47,9\% e 45,4\%, $p = 0{,}5344$). A literatura pós-ChatGPT
não abandona a tese da substituição, mas passa a enfatizar com mais frequência a
complementaridade entre humanos e IA, ao mesmo tempo em que fala menos dos
efeitos de equilíbrio geral. São exemplos da primeira linha, no corpus, a
evidência sobre a origem e o conteúdo de novas ocupações ao longo de quase oito
décadas \parencite{autorNewFrontiers2024} e os ganhos de produtividade da IA
generativa concentrados nos trabalhadores menos experientes
\parencite{brynjolfssonLiRaymond2025}. O fato de a reinstalação não acompanhar a
complementaridade merece registro e é retomado na Discussão: a literatura
recente documenta que a IA \emph{ajuda} o trabalhador nas tarefas existentes com
muito mais frequência do que documenta a criação de tarefas \emph{novas}.

\subsection{Eixo 3: Como se mede}

A forma de medir o efeito também se transforma. O tipo de evidência muda de
maneira marcada ($\chi^2 = 57{,}44$, $p < 0{,}001$, Tabela~\ref{tab:comparacao}).
A mudança mais expressiva é o recuo da modelagem teórica (30,0\% para 13,2\%) e
a multiplicação por quatro dos estudos de firma e de plataformas de
\textit{freelancers} (4,6\% para 19,8\%; p.\,ex.\
\parencite{bessenGoosSalomons2025, hemous2025}). Os estudos de exposição
ocupacional recuam (19,2\% para 14,8\%; típicos do período anterior, p.\,ex.\
\parencite{lordanNeumark2018}), enquanto a evidência macro/setorial (25,4\% e
27,3\%) e as revisões (20,8\% e 24,0\%) mantêm-se aproximadamente estáveis. O
campo desloca-se, assim, do modelo para o dado e do agregado para a firma.

A magnitude dos efeitos, contudo, permanece pouco comparável
(Tabela~\ref{tab:magnitude_cobertura}). Apenas 12,1\% dos estudos do pré e
13,5\% do pós reportam um efeito normalizável; entre esses, a mediana da
magnitude normalizada passa de 0,453 para 0,093, mas as faixas são amplas e as
unidades, heterogêneas, de modo que a diferença não comporta interpretação
substantiva. Coerentemente com o desenho narrativo-estruturado e não
meta-analítico, não se agregam (\textit{pooling}) essas magnitudes nem se aplica
teste de hipótese sobre elas; o sinal qualitativo permanece o eixo principal de
síntese. Essa baixa cobertura, somada ao predomínio de desenhos descritivos
documentado no Capítulo~\ref{cap:descritivas}, é uma limitação substantiva do
estado atual da literatura, não apenas desta revisão.

\begin{table}[htbp]
\centering
\caption{Cobertura e descritivo da magnitude normalizada, pré vs.\ pós-ChatGPT.}
\label{tab:magnitude_cobertura}
\input{tables/magnitude_cobertura.tex}
\end{table}

\subsection{Eixo 4: O que se sabe vs. o que se projeta}

O horizonte dos achados se encurta de forma acentuada após o ChatGPT
($\chi^2 = 66{,}07$, $p < 0{,}001$): os estudos de curto prazo mais que
triplicam de participação (2,9\% para 10,5\%), o médio prazo torna-se
majoritário (40,3\% para 61,8\%), o longo prazo cai quase pela metade (46,6\%
para 24,0\%) e as projeções recuam de 10,1\% para 3,7\%. A literatura recente é,
nesse sentido, menos especulativa e mais ancorada em efeitos observáveis no
presente, como nas primeiras estimativas dos efeitos da IA generativa sobre o
emprego a partir de dados populacionais \parencite{kauhanenRouvinen2025}.

O sinal sobre o emprego também muda ($\chi^2 = 27{,}93$, $p < 0{,}001$), e a
mudança não é a que o debate público sugeriria: o sinal ``negativo'' recua
(30,7\% para 21,6\%) e o ``positivo'' mais que quadruplica (3,8\% para 16,8\%),
enquanto o ``ambíguo'' segue majoritário e praticamente estável (63,4\% e
59,5\%). Em vez de um veredito mais alarmista, a evidência pós-ChatGPT é, no
agregado, \emph{menos pessimista} quanto ao saldo de empregos. O
Capítulo~\ref{cap:discussao} mostra que parte dessa mudança decorre de
\emph{o que} a literatura recente mede, e não apenas do que ela encontra.

\subsection{Continuidades}

Três elementos constituem o núcleo de continuidade entre os dois períodos: (i) o
sinal predominante sobre o emprego é ambíguo antes e depois; (ii) a categoria
modal de polarização continua sendo o risco para a baixa qualificação; e (iii) o
deslocamento permanece o mecanismo mais mobilizado, e a reinstalação é invocada
por pouco menos da metade dos estudos nos dois períodos. A tese da polarização e
a centralidade do efeito-substituição, herdadas da era da automação, sobrevivem
à chegada da IA generativa como moldura dominante.

\subsection{Rupturas}

As rupturas são mais numerosas e mais fortes do que as continuidades, e todas
resistem à correção de Holm para comparações múltiplas
(Tabela~\ref{tab:multiplos_testes}): quintuplica a fração de estudos que situam
o risco na alta qualificação; cresce a complementaridade e recuam o
deslocamento e a demanda agregada como mecanismos; o tipo de evidência migra da
modelagem teórica para os estudos de firma e de plataformas; o horizonte se
encurta; e o sinal torna-se menos negativo e mais frequentemente positivo. O
retrato que emerge não é o de uma inversão do consenso anterior --- a baixa
qualificação segue modal, o sinal segue ambíguo ---, mas tampouco o de mera
mudança de ênfase: é o de um \emph{deslocamento parcial e consistente de
diagnóstico}. A IA generativa amplia o conjunto de ocupações sob escrutínio,
traz o trabalho cognitivo qualificado para dentro da zona de risco e desloca a
discussão da substituição para a complementaridade, sem revogar o diagnóstico de
polarização que a literatura vinha construindo desde a era da automação. O
próximo capítulo examina o que está por trás desse deslocamento.
